\chapter{Appendix}
\label{appendix:deliverables}

\section{Code and Models}
\label{sec:appendix-code}

\begin{enumerate}
  \item Source code (GitHub): \url{https://github.com/magwrap/north-sami-ocr/}
  \item Trained lightweight CNN-CTC model card (HuggingFace): \url{https://huggingface.co/magwrap/cnn-ctc-ocr-sme}
  \item Spr\r{a}kbanken synthetic S\'{a}mi OCR dataset (HuggingFace, mirrored): \url{https://huggingface.co/datasets/magwrap/synthetic_sami_ocr_data}
\end{enumerate}


\section{Conference Submissions}
\label{sec:appendix-submissions}

\begin{enumerate}
  \item \textbf{IISA 2026}, International Conference on Information, Intelligence, Systems and Applications. Full paper accepted to the \textit{21st century AI and education} theme. Venue: \url{https://easyconferences.eu/iisa2026/}.
  \item \textbf{SCAI 2026}, 15th Scandinavian Conference on Artificial Intelligence. Extended abstract accepted to the \textit{AI for Good: Innovation, Ethics, Society} theme. Venue: \url{https://www.sdu.dk/en/forskning/sdu_applied_ai_and_data_science/scai}.
\end{enumerate}

\section{Supervisor and Meeting Log}
\label{sec:appendix-supervision}

Weekly meetings were held on Tuesdays, serving as the primary supervisory mechanism throughout the project. These meetings provided a structured opportunity for regular check ins, progress updates, and feedback sessions. The discussions during these meetings were instrumental in guiding the research direction, troubleshooting issues, and refining methodologies.

\section{AI usage}

This declaration follows SDU's guidelines on generative AI in student work, which require transparent disclosure of how AI contributed to the final product. Throughout the project, generative AI tools, namely Claude Code (Sonnet 4.6, Opus 4.7, Opus 4.8) were utilized for the following categories of tasks:
\begin{itemize}
  \item \textbf{Programming and code generation}: AI assisted with writing code snippets, boilerplate, and data processing scripts. Example: the alignment tool experiment for parallel corpora, which would normally require weeks of development time, was prototyped in hours with AI assistance, freeing time for experimentation and analysis. All AI-generated code was read, tested, and modified before integration.
  \item \textbf{Brainstorming and ideation}: AI was used to generate ideas, suggest alternative approaches, and surface considerations that a less experienced researcher might overlook. All directions explored were subsequently verified against the literature and the supervisor's feedback.
  \item \textbf{Language proofreading}: AI was used to check grammar, spelling, punctuation, and to improve clarity of phrasing. AI-suggested rewordings were reviewed and edited before integration.
  \item \textbf{Structure and organization}: AI was occasionally consulted on chapter and section ordering and on how to frame arguments. 
  \item \textbf{Figure generation}: AI assisted in producing the Python/matplotlib scripts used to generate figures and LaTeX tables (see \texttt{figures/scripts/}). The underlying data and the analytical claims drawn from there were independent of AI.
\end{itemize}

\noindent\textbf{Transparency and accountability.} No AI-generated text or output is presented in this report as the my own work without disclosure. I have critically evaluated every AI contribution for relevance, accuracy, and factual correctness, and take full responsibility for the final content of the report.